\section{Measured Values}
\label{sec:appendix}

This appendix collects the measured values behind the findings, grouped by theme. Every number
here was regenerated from the current code rather than copied from earlier notes. The body of the
paper states what each finding shows. This is the evidence. P-numbers refer to the experiments in
\texttt{code/experiment/} and the result notes in \texttt{note/experiment/results/}.

\subsection{Geometry and the base}

\begin{table}[H]
\centering
\small
\begin{tabular}{lll}
\toprule
\multicolumn{1}{c}{probe scale (momentum)} & \multicolumn{1}{c}{lattice anisotropy} & \multicolumn{1}{c}{sprinkle anisotropy} \\
\midrule
0.5 & 0.016 & $\approx 0.04$ \\
1.0 & 0.067 & $\approx 0.04$ \\
2.0 & 0.332 & $\approx 0.04$ \\
3.0 & 1.121 & $\approx 0.04$ \\
\bottomrule
\end{tabular}
\caption{\label{tab:lorentz}Lorentz isotropy (P27). The lattice's photon-speed anisotropy grows with energy to order one near the cutoff, a real preferred-frame effect. The random sprinkle stays near $0.04$ at every scale, finite-sample noise on an exactly isotropic distribution that falls toward zero as the sample grows.}
\end{table}

\begin{table}[H]
\centering
\small
\begin{tabular}{lllll}
\toprule
\multicolumn{1}{c}{crystal} & \multicolumn{1}{c}{dimension} & \multicolumn{1}{c}{mean degree} & \multicolumn{1}{c}{Lorentz anisotropy} & \multicolumn{1}{c}{safe} \\
\midrule
heptagrid $\{7,3\}$ & 2D & 74.9 & 0.032 & yes \\
pentagrid $\{5,4\}$ & 2D & 60.6 & 0.030 & yes \\
$\{8,3\}$ & 2D & 27.0 & 0.050 & yes \\
$\{6,4\}$ & 2D & 19.1 & 0.039 & yes \\
dodecagrid $\{5,3,4\}$ & 3D & 11.6 & 0.050 & yes \\
\bottomrule
\end{tabular}
\caption{The Coxeter family is Lorentz-safe (P41, P47). Across the regular hyperbolic tilings of the plane and the three-dimensional dodecagrid, the link-direction anisotropy stays small, where a flat lattice reaches order one (Table~\ref{tab:lorentz}).}
\end{table}

\begin{table}[H]
\centering
\small
\begin{tabular}{lll}
\toprule
\multicolumn{1}{c}{space} & \multicolumn{1}{c}{honeycombs} & \multicolumn{1}{c}{examples} \\
\midrule
$H^2$ & infinitely many & all $\{p,q\}$, $1/p+1/q<1/2$ \\
$H^3$ & 4 & $\{3,5,3\}, \{4,3,5\}, \{5,3,4\}, \{5,3,5\}$ \\
$H^4$ & 5 & $\{3,3,3,5\}, \ldots, \{5,3,3,5\}$ \\
$H^5$ & 0 & none \\
$H^6$ and above & 0 & none \\
\bottomrule
\end{tabular}
\caption{The dimension window (P62). Compact regular hyperbolic honeycombs with finite cells exist only in spatial dimensions two, three, and four. The model's substrate is $\{5,3,4\}$ in $H^3$. The testbed reproduces the known classification exactly.}
\end{table}

\begin{table}[H]
\centering
\small
\begin{tabular}{lll}
\toprule
\multicolumn{1}{c}{spatial dimension} & \multicolumn{1}{c}{apsidal angle} & \multicolumn{1}{c}{stable closed orbit} \\
\midrule
2 & $127^\circ$ & no (precesses) \\
3 & $180^\circ$ & yes \\
4 & unstable & no \\
5 & unstable & no \\
\bottomrule
\end{tabular}
\caption{Dimension selection (P68). Within the window, only three spatial dimensions gives stable, closed gravitational orbits. Two precesses (apsidal angle $127^\circ$, not $180^\circ$), and four and above are unstable.}
\end{table}

\begin{table}[H]
\centering
\small
\begin{tabular}{llll}
\toprule
\multicolumn{1}{c}{grown flat mesh} & \multicolumn{1}{c}{target} & \multicolumn{1}{c}{measured} & \multicolumn{1}{c}{fit quality} \\
\midrule
2D & 2 & 2.00 & $r^2 = 1.000$ \\
3D & 3 & 2.97 & $r^2 = 1.000$ \\
4D & 4 & 3.90 & $r^2 = 1.000$ \\
\bottomrule
\end{tabular}
\caption{Emergent dimension from pure growth (P69). The dimension read from a grown flat mesh, with no embedding put in, matches the target. A grown curved mesh instead grows exponentially, the fingerprint of negative curvature.}
\end{table}

\begin{table}[H]
\centering
\small
\begin{tabular}{lll}
\toprule
\multicolumn{1}{c}{mesh size $N$} & \multicolumn{1}{c}{2D estimate} & \multicolumn{1}{c}{3D estimate} \\
\midrule
500 & 1.994 & 2.980 \\
1000 & 2.005 & 2.991 \\
2000 & 2.004 & 2.978 \\
4000 & 2.010 & 2.991 \\
\bottomrule
\end{tabular}
\caption{The continuum limit (P52). The emergent dimension holds at its target across mesh sizes, a stable continuum limit rather than a finite-size artefact.}
\end{table}

\begin{table}[H]
\centering
\small
\begin{tabular}{ll}
\toprule
\multicolumn{1}{c}{elements (after coarse-graining)} & \multicolumn{1}{c}{dimension} \\
\midrule
6000 (original) & 2.005 \\
2990 & 2.016 \\
1499 & 2.018 \\
763 & 2.020 \\
373 & 2.072 \\
\bottomrule
\end{tabular}
\caption{The coarse-graining fixed point (P53). Halving the element count at each level leaves the dimension invariant, a renormalization fixed point (2D sprinkling shown).}
\end{table}

\begin{table}[H]
\centering
\small
\begin{tabular}{ll}
\toprule
\multicolumn{1}{c}{independent construction} & \multicolumn{1}{c}{value} \\
\midrule
continued fraction $x = 1 + 1/x$ & $1.6180340$ \\
ratio of consecutive Fibonacci numbers & $1.6180340$ \\
pentagon geometry $2\cos(\pi/5)$ & $1.6180340$ \\
\bottomrule
\end{tabular}
\caption{Three roads to the golden ratio (P50). The same constant arrives from the crystal's growth law, its Fibonacci addressing, and its pentagonal geometry.}
\end{table}

\begin{table}[H]
\centering
\small
\begin{tabular}{ll}
\toprule
\multicolumn{1}{c}{quantity} & \multicolumn{1}{c}{value} \\
\midrule
measured ball-growth ratio & 2.6186 \\
pentagrid law $(3+\sqrt 5)/2$ & 2.6180 \\
\bottomrule
\end{tabular}
\caption{Deterministic eternal growth (P83). The ball-growth ratio of the one-cell-at-a-time grown crystal converges on its own to the pentagrid's golden-ratio law, so the curvature emerges from the growth.}
\end{table}

\subsection{Dynamics, matter, and gravity}

\begin{table}[H]
\centering
\small
\begin{tabular}{llll}
\toprule
\multicolumn{1}{c}{mass $m$} & \multicolumn{1}{c}{gap $\omega(0)$} & \multicolumn{1}{c}{fitted $a$} & \multicolumn{1}{c}{fitted $b$ ($m^2$)} \\
\midrule
0.0 & 0.000 & 0.968 & 0.000 \\
0.3 & 0.300 & 0.968 & 0.090 \\
0.6 & 0.600 & 0.968 & 0.360 \\
\bottomrule
\end{tabular}
\caption{Mass and the relativistic dispersion (P14). The 1D lattice Dirac operator gives $\omega^2 = a k^2 + b$ with $a$ near one (the light speed) and $b$ near $m^2$ (the rest energy squared). At $m=0$ the gap vanishes.}
\end{table}

\begin{table}[H]
\centering
\small
\begin{tabular}{llllll}
\toprule
\multicolumn{1}{c}{$L$} & \multicolumn{1}{c}{links} & \multicolumn{1}{c}{gauge modes} & \multicolumn{1}{c}{physical modes} & \multicolumn{1}{c}{min $\omega^2$} & \multicolumn{1}{c}{Proca min} \\
\midrule
3 & 81 & 29 & 52 & 3.000 & 1.00 \\
4 & 192 & 66 & 126 & 2.000 & 1.00 \\
5 & 375 & 127 & 248 & 1.382 & 1.00 \\
\bottomrule
\end{tabular}
\caption{The photon (P20). About a third of the link degrees of freedom are exact gauge zero modes, and the smallest physical $\omega^2$ shrinks as the lattice grows, so the photon is massless. A mass term (Proca) would fix a gap near $m^2 = 1$.}
\end{table}

\begin{table}[H]
\centering
\small
\begin{tabular}{ll}
\toprule
\multicolumn{1}{c}{field scale $\epsilon$} & \multicolumn{1}{c}{Wilson action / Maxwell action} \\
\midrule
0.50 & 0.92643 \\
0.20 & 0.98778 \\
0.10 & 0.99693 \\
0.05 & 0.99923 \\
0.02 & 0.99988 \\
\bottomrule
\end{tabular}
\caption{The gauge operator from the action (P23). As the field scale shrinks, the discrete Wilson gauge action converges to the continuum Maxwell action, and the resulting operator is massless.}
\end{table}

\begin{table}[H]
\centering
\small
\begin{tabular}{lll}
\toprule
\multicolumn{1}{c}{quantity} & \multicolumn{1}{c}{model} & \multicolumn{1}{c}{observed} \\
\midrule
$W$ mass (GeV) & 80.0 & 80.4 \\
$Z$ mass (GeV) & 91.3 & 91.2 \\
photon mass (GeV) & 0.000 & 0 \\
$\sin^2\theta_W$ & 0.233 & 0.231 \\
\bottomrule
\end{tabular}
\caption{Electroweak breaking (P25). With the measured couplings and $v = 246$ GeV, the Higgs mechanism reproduces the observed boson masses and Weinberg angle, with the photon left massless.}
\end{table}

\begin{table}[H]
\centering
\small
\begin{tabular}{lll}
\toprule
\multicolumn{1}{c}{dimension} & \multicolumn{1}{c}{form} & \multicolumn{1}{c}{fit quality} \\
\midrule
1D & linear (confining) & $R^2 = 0.94$ \\
2D & logarithmic & $R^2 = 0.98$ \\
3D & Newtonian $1/r$ & $R^2 = 0.9965$ \\
\bottomrule
\end{tabular}
\caption{The Newtonian limit (P16). The gravitational potential takes the right form in each dimension: confining in 1D, logarithmic in 2D, and inverse-distance (Newtonian) in 3D.}
\end{table}

\begin{table}[H]
\centering
\small
\begin{tabular}{ll}
\toprule
\multicolumn{1}{c}{region size $l$} & \multicolumn{1}{c}{entanglement entropy} \\
\midrule
2 & 3.97 \\
3 & 10.67 \\
4 & 21.61 \\
\bottomrule
\end{tabular}
\caption{Black-hole entropy, the area law (P33). The entanglement entropy of an $l \times l \times l$ region scales with area, not volume.}
\end{table}

\begin{table}[H]
\centering
\small
\begin{tabular}{llll}
\toprule
\multicolumn{1}{c}{density $\rho$} & \multicolumn{1}{c}{min length $\ell$} & \multicolumn{1}{c}{curvature cap $1/\ell^2$} & \multicolumn{1}{c}{$\ell\sqrt{\rho}$} \\
\midrule
1 & 0.545 & 3.37 & 0.545 \\
2 & 0.342 & 8.54 & 0.484 \\
4 & 0.230 & 18.95 & 0.459 \\
8 & 0.153 & 42.66 & 0.433 \\
16 & 0.102 & 96.79 & 0.407 \\
\bottomrule
\end{tabular}
\caption{Singularity resolution (P28). The minimum length shrinks with density but the curvature, going as $1/\ell^2$, is capped at a finite value, so the substrate has no infinite-density point.}
\end{table}

\subsection{The quantum sector}

\begin{table}[H]
\centering
\small
\begin{tabular}{lll}
\toprule
\multicolumn{1}{c}{time $t$} & \multicolumn{1}{c}{quantum width} & \multicolumn{1}{c}{classical width} \\
\midrule
4 & 5.66 & 2.83 \\
8 & 11.31 & 4.00 \\
16 & 22.63 & 5.66 \\
32 & 45.25 & 8.00 \\
48 & 67.88 & 9.80 \\
\bottomrule
\end{tabular}
\caption{Quantum coherence (P17). A quantum walk on the mesh spreads ballistically (width $\sim t^{1.0}$), a classical walk only diffusively (width $\sim t^{0.5}$), the signature of coherence.}
\end{table}

\begin{table}[H]
\centering
\small
\begin{tabular}{llll}
\toprule
\multicolumn{1}{c}{branch} & \multicolumn{1}{c}{$|c|^2$ (Born)} & \multicolumn{1}{c}{by counting} & \multicolumn{1}{c}{by envariance} \\
\midrule
1 & 0.040 & 0.040 & 0.040 \\
2 & 0.252 & 0.253 & 0.252 \\
3 & 0.494 & 0.493 & 0.494 \\
4 & 0.213 & 0.213 & 0.213 \\
\bottomrule
\end{tabular}
\caption{The Born rule, derived (P70). Fair counting and Zurek envariance both reproduce $|c|^2$. Fair sampling selects the exponent $p=2$ uniquely (mismatch $0.117$ at $p=1$, $0.001$ at $p=2$, $0.105$ at $p=3$).}
\end{table}

\subsection{Cosmology}

\begin{table}[H]
\centering
\small
\begin{tabular}{ll}
\toprule
\multicolumn{1}{c}{spacetime volume $V$} & \multicolumn{1}{c}{RMS implied $\delta\Lambda$} \\
\midrule
$10^2$ & $9.97\times 10^{-2}$ \\
$10^3$ & $3.16\times 10^{-2}$ \\
$10^4$ & $1.00\times 10^{-2}$ \\
$10^5$ & $3.16\times 10^{-3}$ \\
$10^6$ & $1.00\times 10^{-3}$ \\
$10^7$ & $3.18\times 10^{-4}$ \\
\bottomrule
\end{tabular}
\caption{The everpresent cosmological constant (P46). The fluctuation of $\Lambda$ scales as $V^{-1/2}$ (measured slope $-0.499$), so it shrinks as the universe grows, landing near the observed dark-energy scale at cosmic volume.}
\end{table}

\begin{table}[H]
\centering
\small
\begin{tabular}{ll}
\toprule
\multicolumn{1}{c}{region grid} & \multicolumn{1}{c}{density contrast} \\
\midrule
$4^3$ & 0.017 \\
$6^3$ & 0.032 \\
$8^3$ & 0.052 \\
$12^3$ & 0.091 \\
$16^3$ & 0.144 \\
$24^3$ & 0.262 \\
\bottomrule
\end{tabular}
\caption{The primordial fluctuation seed (P78). The density contrast falls as one over the square root of the count, a scale-free Poisson seed (measured exponent $-0.509$).}
\end{table}

\begin{table}[H]
\centering
\small
\begin{tabular}{llll}
\toprule
\multicolumn{1}{c}{test} & \multicolumn{1}{c}{experimental bound} & \multicolumn{1}{c}{model} & \multicolumn{1}{c}{lattice} \\
\midrule
linear LIV $\xi_1$ & $<0.132$ (Fermi-LAT) & $0$ & $\sim 1.12$ (excluded) \\
quadratic LIV $E_{QG2}$ & $>1.3\times 10^{11}$ GeV & passes & n/a \\
swerve diffusion & cosmic-ray spectra & $\sim$ density$^{-1.61}$ & n/a \\
\bottomrule
\end{tabular}
\caption{Sharp predictions against the bounds (P82). The substrate predicts no first-order Lorentz violation and passes the gamma-ray-burst timing bound that excludes a lattice. The swerve shrinks as the discreteness fines.}
\end{table}

\subsection{The Standard Model content}

\begin{table}[H]
\centering
\small
\begin{tabular}{lll}
\toprule
\multicolumn{1}{c}{field} & \multicolumn{1}{c}{solved hypercharge} & \multicolumn{1}{c}{Standard Model} \\
\midrule
$Q$ (quark doublet) & $0.1667$ & $1/6$ \\
$u$ (up antiquark) & $-0.6667$ & $-2/3$ \\
$d$ (down antiquark) & $0.3333$ & $1/3$ \\
$L$ (lepton doublet) & $-0.5000$ & $-1/2$ \\
$e$ (positron) & $1.0000$ & $1$ \\
\bottomrule
\end{tabular}
\caption{Anomaly cancellation forces charge quantization (P79). The unique anomaly-free, Yukawa-consistent solution for one generation is exactly the Standard Model hypercharges. The cubic and color anomalies then cancel to $\sim 10^{-12}$, and charges come out quantized in thirds.}
\end{table}

\begin{table}[H]
\centering
\small
\begin{tabular}{ll}
\toprule
\multicolumn{1}{c}{setup} & \multicolumn{1}{c}{net asymmetry} \\
\midrule
all three conditions present & 0.081 \\
no C/CP violation & 0.000 \\
in equilibrium & 0.000 \\
no baryon-number violation & 0.000 \\
\bottomrule
\end{tabular}
\caption{Baryogenesis (P80). With all three Sakharov conditions a matter excess builds up, and removing any one erases it.}
\end{table}

\begin{table}[H]
\centering
\small
\begin{tabular}{ll}
\toprule
\multicolumn{1}{c}{overlap law} & \multicolumn{1}{c}{mass spread (decades)} \\
\midrule
exponential (hyperbolic substrate) & 5.5 \\
power-law (flat substrate) & 2.3 \\
observed (top quark over electron) & 5.5 \\
\bottomrule
\end{tabular}
\caption{The mass hierarchy (P81). Exponential overlap on a hyperbolic substrate spans the observed charged-fermion range from evenly spaced modes, where a flat power law cannot.}
\end{table}

\subsection{The interior}

\begin{table}[H]
\centering
\small
\begin{tabular}{ll}
\toprule
\multicolumn{1}{c}{region} & \multicolumn{1}{c}{integration $\Phi$} \\
\midrule
a genuine self (a cell) & 4.397 \\
a random same-size bag & 0.000 \\
the whole mesh & 0.116 \\
\bottomrule
\end{tabular}
\caption{Integrated information (P63). A genuine self resists being cut into parts (high $\Phi$), a random bag of vibes does not, and the whole mesh is one integrated unity. A self sits at a local maximum of $\Phi$.}
\end{table}

\begin{table}[H]
\centering
\small
\begin{tabular}{ll}
\toprule
\multicolumn{1}{c}{coarse-graining} & \multicolumn{1}{c}{macro-rule agreement} \\
\midrule
arbitrary blocks & 0.47 \\
coherent domains (naive) & 0.76 \\
coherent domains (renormalized) & 0.80 \\
domains of size $\ge 3$ & 0.94 \\
domains of size $\ge 5$ & 0.95 \\
domains of size $\ge 10$ & 1.00 \\
\bottomrule
\end{tabular}
\caption{The emergent macro-rule (P58). Coarse-grained along the system's own coherent domains, the higher level obeys the same rule, and the agreement rises with how integrated the domains are. Arbitrary blocks do not work.}
\end{table}

\begin{table}[H]
\centering
\small
\begin{tabular}{lll}
\toprule
\multicolumn{1}{c}{fraction of cell flipped} & \multicolumn{1}{c}{cell returns to body} & \multicolumn{1}{c}{rest of body intact} \\
\midrule
10\% & 1.00 & 1.00 \\
25\% & 1.00 & 1.00 \\
50\% & 0.66 & 1.00 \\
75\% & 0.00 & 1.00 \\
100\% & 0.00 & 1.00 \\
\bottomrule
\end{tabular}
\caption{Nested selves (P59). Flip part of one cell of a body and re-settle. A small wound heals (homeostasis), a whole-cell flip persists with a new identity (autonomy), and the rest of the body is undisturbed throughout.}
\end{table}

\begin{table}[H]
\centering
\small
\begin{tabular}{lllll}
\toprule
\multicolumn{1}{c}{level} & \multicolumn{1}{c}{name} & \multicolumn{1}{c}{units} & \multicolumn{1}{c}{internal alignment} & \multicolumn{1}{c}{rule holds} \\
\midrule
0 & cells & 81 & 0.99 & 1.00 \\
1 & tissues & 27 & 0.84 & 1.00 \\
2 & organs & 9 & 0.59 & 1.00 \\
3 & systems & 3 & 0.48 & 1.00 \\
4 & body & 1 & 0.21 & n/a \\
\bottomrule
\end{tabular}
\caption{The tower of selves (P60). A recursively modular mesh descends to a single top self, dividing by about three at each level, with the emergent rule holding at every level.}
\end{table}

\begin{table}[H]
\centering
\small
\begin{tabular}{ll}
\toprule
\multicolumn{1}{c}{model size (fraction of the whole)} & \multicolumn{1}{c}{reconstruction fidelity} \\
\midrule
2\% & 0.54 \\
5\% & 0.52 \\
10\% & 0.54 \\
25\% & 0.63 \\
50\% & 0.71 \\
100\% & 1.00 \\
\bottomrule
\end{tabular}
\caption{No complete self-storage (P61). A self-model's reconstruction fidelity rises with its size but reaches one only at the full size of the whole, so a lossless self-record cannot fit inside the system.}
\end{table}

\begin{table}[H]
\centering
\small
\begin{tabular}{ll}
\toprule
\multicolumn{1}{c}{perturbation (fraction flipped)} & \multicolumn{1}{c}{recovery} \\
\midrule
0\% & 1.00 \\
10\% & 1.00 \\
20\% & 1.00 \\
30\% & 1.00 \\
40\% & 0.91 \\
50\% & 0.04 \\
\bottomrule
\end{tabular}
\caption{Selves as attractors (P75). A stored self recovers fully from perturbations up to about a forty percent flip, its basin of attraction, and loses its identity only beyond it.}
\end{table}

\subsection{The capstone}

\begin{table}[H]
\centering
\small
\begin{tabular}{ll}
\toprule
\multicolumn{1}{c}{structure read off the one mesh} & \multicolumn{1}{c}{result} \\
\midrule
mean degree (neighbours per vibe) & 10.6 \\
Lorentz anisotropy & 0.070 (no preferred frame) \\
exponential reach (hyperbolic geometry) & yes \\
tones stay strictly ternary & yes \\
dynamics converges to stable states & yes (flip fraction $\to 0.000$) \\
Hamiltonian bounded below & yes (min eigenvalue $0.0010$) \\
arrow of time (relations accumulate) & yes \\
\bottomrule
\end{tabular}
\caption{The capstone (P34). Every structure is read off one committed mesh, not from separate hand-built models, which is the real test of unification.}
\end{table}

